% ==========================================================================
%  IEEE-format paper on MDP-DSL (v2: no em-dashes, with syntax reference,
%  humble tone throughout).
% ==========================================================================
\documentclass[conference]{IEEEtran}
\IEEEoverridecommandlockouts

\usepackage{cite}
\usepackage{amsmath,amssymb,amsfonts}
\usepackage{algorithmic}
\usepackage{graphicx}
\usepackage{textcomp}
\usepackage{xcolor}
\usepackage{listings}
\usepackage{url}
\usepackage[colorlinks=true,linkcolor=blue!50!black,citecolor=blue!50!black,urlcolor=blue!50!black]{hyperref}

\def\BibTeX{{\rm B\kern-.05em{\sc i\kern-.025em b}\kern-.08em
    T\kern-.1667em\lower.7ex\hbox{E}\kern-.125emX}}

\definecolor{codebg}{HTML}{F6F6F6}
\definecolor{codeframe}{HTML}{CCCCCC}

\lstdefinestyle{dsl}{
  basicstyle=\ttfamily\scriptsize,
  backgroundcolor=\color{codebg},
  frame=single,
  rulecolor=\color{codeframe},
  showstringspaces=false,
  breaklines=true,
  columns=fullflexible,
}

\begin{document}

\title{MDP-DSL: A Compiled Domain-Specific Language\\
       for Sequential Decision-Making with\\
       Formal Verification and Failure Diagnosis}

\author{
\IEEEauthorblockN{Charvit Rajani}
\IEEEauthorblockA{Department of Electronics and Communication Engineering\\
Indian Institute of Technology Guwahati\\
Guwahati 781039, Assam, India\\
Roll No. 240102028}
}

\maketitle

% ==========================================================================
\begin{abstract}
We present MDP-DSL, a compiled domain-specific language for Markov
decision processes and their partially observable extensions. The
language is implemented as a single-translation-unit C++17 compiler
with no external dependencies, organized as a five-phase pipeline:
preprocessor, parser, validator with 25 semantic checks, solver, and
verifier. Four solver backends are integrated: value iteration,
policy iteration, Q-learning, and point-based value iteration (PBVI).
A Baum-Welch hidden Markov model fitter bridges learned regime
transitions into MDP dynamics as a first-class language construct.
The compiler exposes a \texttt{VERIFY} block with compile-time
exit-code semantics suitable for continuous integration, and a
six-class autopsy engine that diagnoses reward misspecification
(myopia, hacking, discount cliffs, magnitude imbalance, dead states,
and policy fragility) with single-parameter repair suggestions. A
separate \texttt{mdp\_autopsy} binary consumes execution logs and
reports model-vs-reality divergence with graded verdicts. We
demonstrate the system on the canonical Tiger POMDP, a $4 \times 3$
grid-world, an HMM-bridged portfolio optimizer for Indian equity
markets, and three failure-analysis demonstrations. The full
implementation comprises approximately 13{,}000 lines of code, and
297 test assertions pass under a single \texttt{g++ -O2 -std=c++17}
invocation. The intent of this work is to explore whether a
compiler-first approach to RL specification can usefully complement
the library-based tooling that dominates the current ecosystem.
\end{abstract}

\begin{IEEEkeywords}
Markov decision processes, POMDPs, domain-specific languages,
formal verification, reward specification, hidden Markov models,
point-based value iteration, compiler design.
\end{IEEEkeywords}

% ==========================================================================
\section{Introduction}

Reinforcement-learning toolchains in widespread use, such as
Gymnasium, Stable-Baselines3, and RLlib, optimize the
training-and-deployment loop. They are excellent for prototyping
neural-network policies. They are not designed, however, for what we
argue is the dominant cost in applied RL: the model itself (the
transition kernel, the reward shape, the discount factor, the safety
properties) is buried inside Python class hierarchies and tensor
manipulations. There is no canonical artifact describing what was
assumed, diff-able between commits, formally checkable by tooling.

Three observations motivated this work.

\emph{First, misspecification is an expensive failure mode.} An RL
algorithm that converges to a policy satisfying the reward function
in a way the designer did not intend, that is, specification gaming,
reward hacking, and discount-cliff phase transitions, is folklore in
the community~\cite{krakovna2020specification}. These failures live
in blog posts and conference talks, not in tooling. If the
specification were a language a compiler could read, the compiler
could complain at edit time.

\emph{Second, verification is the natural endpoint of declarativity.}
Hardware description languages have type-checking; database query
languages have query planners. The sequential-decision world has no
comparable artifact: no \texttt{VERIFY} block at the end of the
file. Formal-verification tools for probabilistic systems (PRISM,
STORM) exist but check properties of pre-specified policies; they do
not themselves solve Bellman optimality.

\emph{Third, data and decisions live in different vocabularies.}
HMMs learn latent regime structure from data~\cite{rabiner1989tutorial};
MDPs solve optimal behavior given dynamics~\cite{puterman2014markov}.
A domain expert who wants to write ``fit an HMM on this data file,
then plumb its transition matrix into this MDP'' has no concise way
to express the connection. This work makes that connection a
first-class language affordance.

MDP-DSL is offered as a small contribution toward all three. The
specific deliverables are: a compiled DSL with 25-rule semantic
validation (Section~\ref{sec:lang}); four interchangeable solver
backends including PBVI for POMDPs (Section~\ref{sec:solvers}); a
Baum-Welch HMM fitter bridged into the MDP language at compile time
(Section~\ref{sec:hmm}); and an autopsy engine for six classes of
reward misspecification plus a backwards-solver binary for
model-vs-reality divergence (Section~\ref{sec:autopsy}). We validate
the system on canonical benchmarks and an HMM-bridged portfolio
optimizer (Sections~\ref{sec:experiments}).

% ==========================================================================
\section{Language Reference}\label{sec:lang}

This section is a complete reference for the surface syntax. A
program is a sequence of declarations, one per line, with a leading
keyword in upper case. Tokens are whitespace-separated; the
\texttt{\#} character introduces a line comment. Keywords are
case-insensitive. No nested scopes or expressions are supported by
design; users who want arithmetic in reward declarations are
expected to precompute the values.

\subsection{Core MDP declarations}

\begin{lstlisting}[style=dsl]
STATE: <name>
ACTION: <name>
TRANSITION: <state> <action> <state> <probability>
REWARD: <state> <value>
ACTION_REWARD: <state> <action> <value>
DISCOUNT: <value>           # gamma in [0, 1)
SOLVER: vi | pi | ql | pbvi
\end{lstlisting}

Probabilities must lie in $[0, 1]$. For every $(s, a)$ pair, the
declared transitions must sum to $1.0$ within $10^{-9}$. Either
\texttt{REWARD} or \texttt{ACTION\_REWARD} may be used (the latter
generalizes the former); they may be mixed.

\subsection{POMDP extensions}

\begin{lstlisting}[style=dsl]
OBSERVATION: <name>
OBSERVE_PROB: <state> <action> <obs> <probability>
INITIAL_BELIEF: <state> <probability>
BELIEF_STATE: <name> <state>:<prob> <state>:<prob> ...
ALPHA_VECTORS:   <int>    # PBVI parameter (default 15)
PBVI_ITERATIONS: <int>    # default 100
HORIZON: <int>            # finite-horizon POMDP
\end{lstlisting}

When \texttt{SOLVER: pbvi}, at least one
\texttt{INITIAL\_BELIEF} or \texttt{BELIEF\_STATE} must be
declared. \texttt{BELIEF\_STATE} entries provide named beliefs that
can be referenced in \texttt{VERIFY} assertions.

\subsection{HMM bridge}

\begin{lstlisting}[style=dsl]
HMM_STATES: <name> <name> ...
EMISSION_TYPE: gaussian | discrete
FIT_HMM: <csv_path> iterations=<int>
HMM_SEED: <int>           # for reproducibility
BRIDGE: hmm -> mdp
\end{lstlisting}

The bridge runs Baum-Welch on the referenced CSV before validation;
the learned transition matrix is plumbed into the AST as
\texttt{TRANSITION:} rows, and the HMM states are added to the
declared state set. User-specified transitions for the same triples
are preserved.

\subsection{Verification}

\begin{lstlisting}[style=dsl]
VERIFY: pi(<state_or_belief>) == <action>
VERIFY: pi(<state_or_belief>) != <action>
VERIFY: V(<state_or_belief>) >  <number>
VERIFY: V(<state_or_belief>) <  <number>
VERIFY: V(<state_or_belief>) >  V(<state_or_belief>)
\end{lstlisting}

Both ASCII (\texttt{pi(s)}) and Unicode ($\pi(s)$) forms are
accepted. The compiler exits with status 1 if any assertion fails,
0 otherwise, which is intended to fit into CI pipelines.

\subsection{Macros}

\begin{lstlisting}[style=dsl]
GRID: <rows>x<cols>:<spec_strings>
\end{lstlisting}

The \texttt{GRID:} macro expands into the equivalent STATE,
TRANSITION, and REWARD declarations for a rectangular grid world,
with optional walls, terminal cells, and slip probability.

\subsection{Complete example: Tiger POMDP}

Listing~\ref{lst:tiger} is the canonical POMDP benchmark expressed
in the language. The compiler solves this with PBVI in 100 backups,
producing 4 useful alpha-vectors after witness pruning, and both
\texttt{VERIFY} assertions pass at exit code 0.

\begin{lstlisting}[style=dsl,caption={Tiger POMDP in MDP-DSL.},label={lst:tiger}]
STATE: TigerLeft
STATE: TigerRight

ACTION: Listen
ACTION: OpenLeft
ACTION: OpenRight

TRANSITION: TigerLeft  Listen TigerLeft  1.0
TRANSITION: TigerRight Listen TigerRight 1.0
# (reset transitions for open actions omitted)

OBSERVATION: HearLeft
OBSERVATION: HearRight

OBSERVE_PROB: TigerLeft  Listen HearLeft  0.85
OBSERVE_PROB: TigerLeft  Listen HearRight 0.15
OBSERVE_PROB: TigerRight Listen HearLeft  0.15
OBSERVE_PROB: TigerRight Listen HearRight 0.85

ACTION_REWARD: TigerLeft  OpenLeft -100.0
ACTION_REWARD: TigerRight OpenLeft   10.0
ACTION_REWARD: TigerLeft  OpenRight  10.0
ACTION_REWARD: TigerRight OpenRight -100.0
ACTION_REWARD: TigerLeft  Listen     -1.0
ACTION_REWARD: TigerRight Listen     -1.0

BELIEF_STATE: Uniform        TigerLeft:0.5  TigerRight:0.5
BELIEF_STATE: AlmostSureLeft TigerLeft:0.95 TigerRight:0.05

DISCOUNT: 0.95
SOLVER:   pbvi

VERIFY: pi(Uniform)        == Listen
VERIFY: pi(AlmostSureLeft) == OpenLeft
\end{lstlisting}

% ==========================================================================
\section{Compilation Pipeline}

The compiler processes input through five mandatory phases, plus an
optional HMM-bridge phase if requested.

\emph{Phase 0, Preprocessor.} Expands \texttt{GRID:} macros into
ordinary STATE / TRANSITION / REWARD lines, which the subsequent
phases process uniformly.

\emph{Phase 1, Parser.} A hand-written line-and-keyword recursive
descent reads each declaration and populates a flat
\texttt{MDP\_AST} struct. The grammar is regular at file level and
LL(1) per line; no separate tokenizer is required. Errors are
reported with line numbers and stop the pipeline.

\emph{HMM bridge phase.} If the file declares \texttt{BRIDGE: hmm
-> mdp}, the bridge runs Baum-Welch on the referenced CSV, injects
hidden states as MDP states, and adds the learned transition matrix
as \texttt{TRANSITION:} rows. User-specified transitions for the
same triples are preserved.

\emph{Phase 2, Validator.} 25 named semantic checks read the AST
and accumulate errors. Examples include: $\sum_{s'} T(s, a, s') = 1$
for every $(s, a)$ pair within $10^{-9}$ (Check 5);
$\sum_o O(s', a, o) = 1$ for every $(s', a)$ (Check 15);
$\sum_s b_0(s) = 1$ for the initial belief (Check 17). Each check
has positive and negative unit tests in the suite.

\emph{Phase 3, Solver.} The dispatcher consults
\texttt{ast.solver\_mode} and invokes one of four backends
(Section~\ref{sec:solvers}), all returning a uniform
\texttt{SolverResult}.

\emph{Phase 4, Verifier.} Each \texttt{VERIFY:} assertion is parsed
and evaluated. The process exits with status 1 if any assertion
fails, 0 otherwise.

% ==========================================================================
\section{Solver Backends}\label{sec:solvers}

\subsection{Value iteration}

We iterate the Bellman optimality operator
\begin{equation}
(T^* V)(s) = \max_a \Big[ R(s,a) + \gamma \sum_{s'} T(s,a,s') V(s') \Big]
\label{eq:bellman-opt}
\end{equation}
from $V_0 \equiv 0$ until $\|V_{k+1} - V_k\|_\infty < \theta = 10^{-9}$.
$T^*$ is a $\gamma$-contraction in sup-norm, so by the Banach
fixed-point theorem the iteration converges geometrically with rate
$\gamma$. Per-iteration cost is $O(|T|)$ where $|T|$ is the number
of declared transition rows; iteration count is
$O(\log\theta / \log\gamma)$.

\subsection{Policy iteration}

PI alternates two steps. \emph{Evaluation:} solve the linear system
$(I - \gamma P^\pi) V^\pi = R^\pi$ where $P^\pi$ is the
$|\mathcal{S}| \times |\mathcal{S}|$ matrix of transitions under
$\pi$. \emph{Improvement:} compute the greedy policy with respect to
$V^\pi$. Termination is detected when the policy stops changing.

The evaluation step uses Gaussian elimination with partial pivoting.
Without pivoting, ill-conditioned matrices (small self-transition
probabilities, or $\gamma$ approaching 1) produced numerically
meaningless solutions during development. The pivoting addition was
approximately fifteen lines of code and eliminated a class of bugs.

\subsection{Q-learning}

Model-free Q-learning with $\epsilon$-greedy exploration is included
to verify that the same AST is solvable by both model-based and
model-free methods. The step-size schedule
$\alpha_t = 1 / (1 + N(s, a))$ satisfies the Robbins-Monro
conditions; $\epsilon$ decays linearly from 1.0 to 0.05 over 25{,}000
steps. On the $4 \times 3$ grid-world after 50{,}000 episodes, the
recovered policy agrees with $\pi^*$ on at least 8 of 11 states.

\subsection{Point-based value iteration}

For POMDPs, we implement PBVI~\cite{pineau2003point}. The value
function on $\Delta(\mathcal{S})$ is represented as the upper
envelope of a set $\Gamma$ of alpha-vectors. The backup at a belief
point $b$ proceeds in three stages:

\begin{enumerate}
\item For each action $a$ and observation $o$, find the previous
alpha-vector to follow:
\[
\alpha^*_{a,o} = \arg\max_{\alpha \in \Gamma_\text{prev}} \sum_s b(s) \sum_{s'} T(s,a,s') O(s',a,o) \alpha(s').
\]
\item Construct the action-conditional vector:
\[
\alpha_a(s) = R(s,a) + \gamma \sum_o \sum_{s'} T(s,a,s') O(s',a,o) \alpha^*_{a,o}(s').
\]
\item Return $\alpha_{a^*}$ where $a^* = \arg\max_a \langle b, \alpha_a \rangle$.
\end{enumerate}

\emph{Witness pruning.} We retain $\alpha$ if and only if there
exists a belief $b$ in the current belief set $B$ at which $\alpha$
is strictly best. An earlier implementation used general dominance
(that is, $\alpha$ dominated iff some $\alpha'$ exceeds $\alpha$ on
every state) and over-pruned: the Tiger problem collapsed to a single
alpha-vector and produced an incorrect policy at the uniform belief.
Witness pruning yields the expected 3 to 5 useful vectors and the
correct policy. The fix is recorded as a regression test that asserts
$|\Gamma| \in [3, 7]$ on Tiger.

% ==========================================================================
\section{HMM-MDP Bridging}\label{sec:hmm}

The Baum-Welch fitter is implemented in pure C++17 with Rabiner's
scaling discipline. For each timestep, the scaling factor
$c_t = 1 / \sum_i \alpha_t(i)$ normalizes the forward variables; the
log-likelihood is recovered as
$\log P(O \mid \lambda) = -\sum_t \log c_t$. Without scaling,
forward variables underflow to zero for sequences beyond
approximately 50 steps with Gaussian emissions of moderate variance,
which breaks the M-step. A dedicated test verifies the fitter on a
500-step synthetic sequence.

Listing~\ref{lst:portfolio} shows the portfolio model. The user
declares four hidden states and points the bridge at a CSV; the
fitted transition matrix is plumbed into the MDP and the file is
validated and solved as a single artifact.

\begin{lstlisting}[style=dsl,caption={HMM-bridged portfolio optimizer (excerpt).},label={lst:portfolio}]
HMM_STATES:    Bull Bear Sideways Crisis
EMISSION_TYPE: gaussian
FIT_HMM:       portfolio_returns.csv iterations=200
HMM_SEED:      42

BRIDGE: hmm -> mdp

ACTION: AggressiveEquity
ACTION: ModerateEquity
ACTION: Defensive
ACTION: CashAndGold

# Sharpe-calibrated rewards (excerpt):
ACTION_REWARD: Bull   AggressiveEquity  0.85
ACTION_REWARD: Bear   Defensive         0.10
ACTION_REWARD: Crisis CashAndGold       0.08

DISCOUNT: 0.92
SOLVER:   pi

VERIFY: pi(Bull)   == AggressiveEquity
VERIFY: pi(Bear)   == Defensive
VERIFY: pi(Crisis) == CashAndGold
VERIFY: V(Bull)    >  V(Bear)
VERIFY: V(Crisis)  <  V(Sideways)
\end{lstlisting}

The portfolio sub-system fetches NSE/BSE equity prices via
\texttt{yfinance} (15-minute local cache, seed-price fallback when
offline), fits the HMM on monthly returns, calibrates per-regime
Sharpe ratios from Viterbi-decoded segments, emits the
\texttt{.mdp}, invokes the compiler with \texttt{-{}-json}, and
proposes a portfolio allocation. A dashboard, implemented as a
single HTML file with no npm dependencies, renders the regime
history, the policy boundary heat-map over belief space, the
allocation donut, and the \texttt{VERIFY} status panel.

% ==========================================================================
\section{The Autopsy Engine}\label{sec:autopsy}

The autopsy engine diagnoses six classes of reward misspecification
statically, before or independently of any training process.

\subsection{Six failure classes}

\emph{Reward myopia.} The effective planning horizon
$H_\text{eff} = \log\varepsilon / \log\gamma$ with
$\varepsilon = 10^{-3}$. BFS computes the minimum distance
$d_\text{goal}$ from each state to any positive-reward state. If
$d_\text{goal} > H_\text{eff}$, the agent is structurally blind to
the goal. The recommended fix is the minimum
$\gamma' = \exp(\log\varepsilon / d_\text{goal})$.

\emph{Reward hacking.} The optimal-policy graph (each state's
outgoing edge follows $\pi^*$, mode-collapsed) is analyzed via
Tarjan's strongly-connected-components algorithm. For each
non-singleton SCC $C$ with reward sum $r_C$, the cycle value under
optimal play is $V_\text{cycle} = r_C / (1 - \gamma^{|C|})$. The
detector fires if $V_\text{cycle} > 1.05 \cdot V^*$ at the best
state outside $C$. Absorbing terminal goals are excluded from cycle
analysis.

\emph{Discount cliff.} The optimal policy is computed at 200 points
across $\gamma \in [0.01, 0.999]$; cliffs are values at which
$\pi(s)$ flips between adjacent grid points for any state $s$. The
detector reports the cliff nearest the user's $\gamma$ along with
the precise flip.

\emph{Magnitude imbalance.} For each non-zero reward term, the
contribution $c_i = |r_i| / \sum_j |r_j|$ is computed. Terms below
$0.001$ are flagged as effectively ignored; the recommended
scale-up is the minimum to reach 5\% contribution.

\emph{Dead-state detection.} Reverse BFS from the set
$\{s : R(s) > 0 \lor \exists a.\ R(s, a) > 0\}$ identifies states
that no policy can reach a reward from.

\emph{Fragility analysis.} For each $(s, \pi^*(s))$ pair and each
perturbable parameter (every reward and $\gamma$), binary search
finds the smallest $|\delta|$ that flips $\pi^*(s)$. The report
ranks the top-5 most fragile decisions and computes a robustness
score.

\subsection{REPAIR command}

Given a target assertion such as \texttt{pi(Start) == MoveRight},
the \texttt{-{}-repair} mode parses the assertion, checks if it
already holds, and otherwise binary-searches each candidate
parameter for the smallest perturbation that satisfies it. The
proposals are returned ranked by magnitude. Multi-parameter joint
repair is noted as future work; the current implementation handles
only single-parameter searches.

\subsection{Backwards solver}

A separate binary, \texttt{mdp\_autopsy}, consumes an execution log
ending in a \texttt{FAILURE} marker and produces a graded verdict.
Empirical transition frequencies $\hat T(s, a, s')$ are computed,
skipping transitions whose source is an absorbing state in the
model (this filters mission-reset events such as physical reset
after collision). The KL divergence between empirical and model
transitions, with additive smoothing $\varepsilon = 10^{-10}$,
ranks the suspect pairs.

The backwards solver itself answers: what is the smallest
perturbation $\Delta T$ along the empirical direction at which a
\texttt{VERIFY} assertion fails? We parameterize
$T_\alpha = (1 - \alpha) T_\text{model} + \alpha \hat T$ and
binary-search $\alpha \in [0, 1]$ for the smallest value at which
any assertion fails.

\emph{Scope statement.} A free $L^1$ minimization over the
constrained simplex would require a QP solver, which would violate
our zero-dependency constraint. Linear search along the empirical
direction is an honest approximation: the empirical data is the
direction in which the world disagrees with the model, which is
exactly the axis the autopsy is meant to investigate. This
limitation is reported in the autopsy output verbatim.

The verdict ladder is shown in Table~\ref{tab:verdict}.

\begin{table}[h]
\caption{Verdict ladder for \texttt{mdp\_autopsy}.}
\label{tab:verdict}
\centering
\begin{tabular}{lcr}
\hline
\textbf{Worst factor} & \textbf{Verdict} & \textbf{Exit code} \\
\hline
$< 2 \times$            & \texttt{MODEL\_ACCURATE}   & 0 \\
$2 \times$ to $10 \times$ & \texttt{MODEL\_OPTIMISTIC} & 1 \\
$\ge 10 \times$           & \texttt{MODEL\_DANGEROUS}  & 2 \\
\hline
\end{tabular}
\end{table}

% ==========================================================================
\section{Experimental Results}\label{sec:experiments}

\subsection{Test suite}

The full test suite comprises 174 tests and 297 assertions,
distributed as in Table~\ref{tab:tests}. The complete suite executes
in roughly one second.

\begin{table}[h]
\caption{Test suite distribution across project phases.}
\label{tab:tests}
\centering
\begin{tabular}{lrr}
\hline
\textbf{Group} & \textbf{Tests} & \textbf{Assertions} \\
\hline
Baseline (parser / validator / VI / PI / QL) & 94 & 142 \\
POMDP / PBVI / belief update                  & 21 & 47 \\
HMM / Baum-Welch / bridge                     & 15 & 36 \\
Autopsy engine (6 classes plus REPAIR)        & 28 & 43 \\
Backwards solver / log parser                 & 16 & 29 \\
\hline
\textbf{Total}                                & \textbf{174} & \textbf{297} \\
\hline
\end{tabular}
\end{table}

\subsection{Tiger POMDP}

On Tiger ($|\mathcal{S}| = 2$, $|\mathcal{A}| = 3$, $|\Omega| = 2$,
$\gamma = 0.95$), PBVI converges in roughly 100 backups, yielding 4
useful alpha-vectors after witness pruning. Both \texttt{VERIFY}
assertions pass.

\subsection{Portfolio model}

On the HMM-bridged portfolio model (4 hidden states, 4 actions, 60
months of Nifty returns, $\gamma = 0.92$, PI solver), all five
\texttt{VERIFY} assertions pass. The fitted regime persistence
matrix has diagonal entries $\{0.83, 0.79, 0.71, 0.62\}$, broadly
consistent with macroscopic Indian-market behavior over the training
window.

\subsection{Autopsy detections}

On crafted demonstrations:
\begin{itemize}
\item Myopic corridor (9 states, $\gamma = 0.30$, asymmetric action
cost): RewardMyopia ERROR with $H_\text{eff} = 5.7$ and
$d_\text{goal} = 8$; DiscountCliff WARNING at $\gamma^* = 0.27$.
\item Reward hacking demo (cycle $S_2 \leftrightarrow S_3$ yielding
$+12$/step, one-shot goal $+50$): RewardHacking ERROR with
$V_\text{cycle} = 246.15$ exceeding $V^*$ outside the cycle, which
was 228.00.
\item Dead-state demo (disconnected trap cycle): DeadState ERROR
with both trap states identified.
\item Repair demo (3-state corridor, $\gamma = 0.15$, move cost
$-5$): the repair search finds 6 fixes; the minimum is
$\gamma$ from 0.15 to 0.218, that is, $\Delta = 0.068$.
\end{itemize}

\subsection{Backwards-solver verdicts}

Table~\ref{tab:autopsy} summarizes four demonstration logs.

\begin{table}[h]
\caption{\texttt{mdp\_autopsy} verdicts on four demonstrations.}
\label{tab:autopsy}
\centering
\begin{tabular}{lll}
\hline
\textbf{Log} & \textbf{Verdict} & \textbf{Backwards solver} \\
\hline
medical\_clean.log  & \texttt{ACCURATE}   & no failure to back-solve \\
medical\_run.log    & \texttt{OPTIMISTIC} & $\alpha{=}0.04$, $\|\Delta T\|_1{=}0.04$ \\
robot\_run.log      & \texttt{DANGEROUS}  & VERIFY robust under empirical \\
portfolio\_run.log  & \texttt{DANGEROUS}  & $\alpha{=}0.32$, $\|\Delta T\|_1{=}0.76$ \\
\hline
\end{tabular}
\end{table}

The robot demonstration is worth a comment. The empirical
$(\text{NearObstacle}, \text{Advance}) \to \text{Collision}$
transition occurred at 0.25 vs.\ the model's 0.02, a factor of
12.5x. Yet the backwards solver reports that no perturbation along
the empirical direction breaks any \texttt{VERIFY} assertion. Under
empirical transitions, the optimal action at NearObstacle is
SlowDown, which is actually safer in reality than in the model. The
diagnostic message is therefore that the engineer's safety
properties were correctly designed, and the run-time failure
occurred because the deployed policy ignored the model and selected
Advance repeatedly, not because the model's safety specification
was inadequate. We find this kind of distinction (design-level
vs.\ deployment-level failure) useful and not always easy to make
by hand.

% ==========================================================================
\section{Related Work}

\emph{POMDP solvers.} APPL/SARSOP~\cite{kurniawati2008sarsop} and
HSVI~\cite{smith2004heuristic} implement modern point-based solvers
focused on performance. The PBVI implementation in MDP-DSL is a
teaching-grade implementation appropriate to the project's scope and
not intended to compete on solver throughput.

\emph{Probabilistic verification.} PRISM~\cite{kwiatkowska2011prism}
and STORM~\cite{dehnert2017storm} verify temporal-logic properties
of Markov chains and MDPs. They accept rich property languages but
require the policy to be specified or enumerated; they do not
themselves solve Bellman optimality. MDP-DSL goes the other way:
solve the MDP, then verify properties of the optimal policy.

\emph{RL specification gaming.} Krakovna et
al.~\cite{krakovna2020specification} catalog 60+ examples of agents
exploiting reward specifications in unintended ways. The autopsy
engine here is a small tooling response: several of these failures
are detectable statically from the reward function alone (cycles,
dead states, discount cliffs).

\emph{DSLs for sequential decisions.} POMDP file formats (such as
Cassandra's .pomdp format) exist as serialization formats for
solvers, but typically lack semantic checking, verification, or HMM
bridging. MDP-DSL adds these in a single artifact.

% ==========================================================================
\section{Engineering Constraints and Honest Scope}

The implementation enforces five constraints. First, no external C++
dependencies (STL only). Second, single-translation-unit compilation
for the main binary. Third, byte-identical backward compatibility
across all language extensions. Fourth, exit-code semantics suitable
for CI. Fifth, total source under roughly 13{,}000 lines.

Three scope limitations are documented explicitly in the
implementation and report output.

\emph{Backwards solver approximation.} The Phase 3B backwards solver
searches along the empirical direction rather than freely minimizing
$\|\Delta T\|_1$ over the constrained simplex (which would require a
QP solver). The implementation reports this approximation
explicitly.

\emph{REPAIR single-parameter only.} The repair search perturbs one
parameter at a time. Multi-parameter joint repair is future work.

\emph{Underestimation factor cap.} When the model assigns
$p < 10^{-6}$ to an empirically-observed transition, the raw ratio
is mathematically unbounded. We cap the displayed factor at
$1000\times$ to keep reports informative; the KL ranking remains
well-defined via additive smoothing.

We argue that explicit scope statements of this kind are more
credible in tooling intended for verification and diagnosis than
implicit claims of optimality.

% ==========================================================================
\section{Author's Contribution Statement}

The project was specified, architected, and led by the author. The
v2.0 baseline (value iteration, policy iteration, Q-learning, the
\texttt{VERIFY} block, the GRID macro, and the original five-phase
pipeline) was completed in the spring 2026 term as the author's
DA 221M Artificial Intelligence minor-project term work under the
supervision of Prof.\ Shyamanta M.\ Hazarika at IIT Guwahati.

The v3.0 extensions described in this paper (the POMDP/PBVI
extension, the HMM-MDP bridge, the portfolio sub-system, the
six-class autopsy engine, and the backwards-solver binary) are
post-coursework work undertaken by the author independent of the
supervised project. Implementation of these extensions used
AI-assisted code generation under the author's direction. Every
architectural choice, scope constraint, algorithm selection, and
verification semantic was authored by the candidate. AI assistance
was used at the implementation level. This is disclosed for
accuracy.

% ==========================================================================
\section{Conclusion and Future Work}

MDP-DSL is a small contribution toward a compiler-first approach to
sequential decision-making. The 25-rule validator catches model
misspecification before any solver runs; the \texttt{VERIFY} block
turns model assumptions into compile-time assertions; the autopsy
engine diagnoses six classes of reward pathology and proposes
repairs.

The longer-term direction the author intends to pursue is to grow
the project from a compiler into an integrated development
environment for sequential-decision specification: a workspace where
practitioners can write a model, see solver and autopsy output
interactively, iterate on the specification, and version it as a
first-class artifact. Smaller future directions include inverse MDP
compilation (algebraic recovery of MDP structure from observed
policies), discount-cliff cartography, and extending the autopsy
machinery beyond MDPs to other forms of objective specification.

The implementation is open-source under the MIT license. The build
is reproducible with a single \texttt{g++} invocation and no
external dependencies. Issues, bug reports, and contributions are
welcome.

\section*{Acknowledgments}

The author thanks Prof.\ Shyamanta M.\ Hazarika (IIT Guwahati) for
supervising the v2.0 baseline as part of the DA 221M course.

% ==========================================================================
\begin{thebibliography}{99}

\bibitem{krakovna2020specification}
V.~Krakovna et al., ``Specification gaming examples in AI,''
DeepMind Safety Research Blog, 2020.

\bibitem{puterman2014markov}
M.~L. Puterman, \emph{Markov Decision Processes: Discrete
Stochastic Dynamic Programming}, John Wiley \& Sons, 2014.

\bibitem{rabiner1989tutorial}
L.~R. Rabiner, ``A tutorial on hidden Markov models and selected
applications in speech recognition,'' \emph{Proc.\ IEEE}, vol.\ 77,
no.\ 2, pp.\ 257--286, 1989.

\bibitem{pineau2003point}
J.~Pineau, G.~Gordon, and S.~Thrun, ``Point-based value iteration:
An anytime algorithm for POMDPs,'' in \emph{IJCAI}, 2003,
pp.\ 1025--1032.

\bibitem{smith2004heuristic}
T.~Smith and R.~Simmons, ``Heuristic search value iteration for
POMDPs,'' in \emph{UAI}, 2004, pp.\ 520--527.

\bibitem{kurniawati2008sarsop}
H.~Kurniawati, D.~Hsu, and W.~S. Lee, ``SARSOP: Efficient
point-based POMDP planning by approximating optimally reachable
belief spaces,'' in \emph{Robotics: Science and Systems}, 2008.

\bibitem{kwiatkowska2011prism}
M.~Kwiatkowska, G.~Norman, and D.~Parker, ``PRISM 4.0: Verification
of probabilistic real-time systems,'' in \emph{CAV}, 2011,
pp.\ 585--591.

\bibitem{dehnert2017storm}
C.~Dehnert, S.~Junges, J.-P. Katoen, and M.~Volk, ``A storm is
coming: A modern probabilistic model checker,'' in \emph{CAV},
2017, pp.\ 592--600.

\bibitem{sutton2018reinforcement}
R.~S. Sutton and A.~G. Barto, \emph{Reinforcement Learning: An
Introduction}, 2nd ed., MIT Press, 2018.

\bibitem{sondik1971optimal}
E.~J. Sondik, ``The optimal control of partially observable Markov
processes,'' Ph.D. dissertation, Stanford University, 1971.

\end{thebibliography}

\end{document}
